\Article{10}
Third, let the surface of a sphere with radius \(=a\) be represented in a plane. Here we set
\[
\begin{aligned}
x&=a\cos t\mathord{.}\sin u\\
y&=a\sin t\mathord{.}\sin u\\
z&=a\cos u
\end{aligned}
\]
\noindent which gives
\[
\omega=aa\sin u^2\,\mathrm{d}t^2+aa\,\mathrm{d}u^2
\]
\noindent The differential expression \(\omega=0\) therefore gives
\[
\mathrm{d}t\mp i\mathord{.}\frac{\mathrm{d}u}{\sin u}=0
\]
\noindent and integrating it gives
\[
t\pm i\log\mathrm{cotang}\tfrac{1}{2}u=\text{Const.}
\]
\par\indent\SourceLeaf{209}%
Thus, if we again use the symbol \(f\) to denote an arbitrary function, \(X\) must be set equal to the real part and \(iY\) to the imaginary part of
\[
f(t+i\log\mathrm{cotang}\tfrac{1}{2}u)
\]
\noindent We will give a few special cases of this general solution.

If we choose a linear function for \(f\) by setting \(f\upsilon=k\upsilon\), we get
\[
X=kt,\quad Y=k\log\mathrm{cotang}\tfrac{1}{2}u
\]
\noindent Applied to the Earth, this is evidently the same as \textsc{Mercator}'s projection if \(t\) denotes geographic longitude and \(90^0-u\) latitude. The formulas of Article 7 give the ratio of magnification as
\[
m=\frac{k}{a\sin u}
\]

If we take an imaginary exponential function for \(f\), beginning with the simplest one \(f\upsilon=ke^{i\upsilon}\), we get
\[
f(t+i\log\mathrm{cotang}\tfrac{1}{2}u)=ke^{\log\mathrm{tang}\frac{1}{2}u+it}=k\,\mathrm{tang}\tfrac{1}{2}u(\cos t+i\sin t)
\]
\noindent and
\[
X=k\,\mathrm{tang}\tfrac{1}{2}u\mathord{.}\cos t,\quad
Y=k\,\mathrm{tang}\tfrac{1}{2}u\mathord{.}\sin t
\]
\noindent which, as is easy to see, is the polar stereographic projection.

More generally, if we set \(f\upsilon=ke^{i\lambda\upsilon}\), we get
\[
X=k\,\mathrm{tang}\tfrac{1}{2}u^{\lambda}\mathord{.}\cos\lambda t,\quad
Y=k\,\mathrm{tang}\tfrac{1}{2}u^{\lambda}\mathord{.}\sin\lambda t
\]
\noindent For the ratio of magnification we obtain
\[
n=aa\sin u^2,\quad N=1,\quad\varphi\upsilon=i\lambda ke^{i\lambda\upsilon}
\]
\noindent and hence
\[
m=\frac{\lambda k\,\mathrm{tang}\frac{1}{2}u^{\lambda}}{a\sin u}
\]

We see here that the representation of all points for which \(u\) is constant falls on a single circle, and that of all points for which \(t\) is constant falls on a single straight line. We also see that the circles belonging to the different values of \(u\) are concentric. This gives a very useful map projection when only a part of the spherical surface is to be represented. In that case it is best to choose \(\lambda\) so that the ratio of magnification is the same at the
\SourceLeaf{210}
extreme values of \(u\), and thus reaches its smallest value toward the middle. If these extreme values of \(u\) are \(u^0\) and \(u'\), we must therefore set
\[
\lambda=\frac{\log\sin u'-\log\sin u^0}{\log\mathrm{tang}\frac{1}{2}u'-\log\mathrm{tang}\frac{1}{2}u^0}
\]
\noindent Sheets Nos. 19---26 of Professor \textsc{Harding}'s star charts were drawn using this projection.

\Article{11}
The general solution for the example treated in the previous article can also be given in another form. We think we should include it here for its elegance.

It follows from what was set out in Article 6 that, since
\[
\mathrm{tang}\tfrac{1}{2}u(\cos t+i\sin t)
\]
\noindent is a function of
\[
t+i\log\mathrm{cotang}\tfrac{1}{2}u
\]
\noindent and
\[
\mathrm{tang}\tfrac{1}{2}u(\cos t+i\sin t)=\frac{\sin u\cos t+i\sin u\sin t}{1+\cos u}=\frac{x+iy}{a+z}
\]
\noindent the general solution can also be represented by
\[
X+iY=f\frac{x+iy}{a+z},\quad X-iY=f'\frac{x-iy}{a+z}
\]
\noindent That is, \(X\) must be set equal to the real part and \(iY\) to the imaginary part of \(f\frac{x+iy}{a+z}\), where \(f\) denotes an arbitrary function. In place of \(f\frac{x+iy}{a+z}\) we may also take, as is easy to see, an arbitrary function of \(\frac{y+iz}{a+x}\) or of \(\frac{z+ix}{a+y}\).

